\scanpage{67}
\mspage{21}
\[
(23)\qquad I\xrightarrow{\underline{d}}\mathbb Z
\qquad (\text{i.e } I=\coprod_m I_m)
\]
playing the role of a dimension function
(one would have a presentation in some sense
dual with a codimension-type function),
satisfying the condition
\[
(24)\qquad\text{If }i,j\in I\,\deleted{\unsure{\varnothing}},\ i\leq j,\text{ then }d(j)=d(i)+1\Longleftrightarrow j\text{ successor of }i
\]
Let us then set, for $m,n_{\deleted{\unsure{2|q}}}\in\mathbb Z$
\[
(25)\qquad\left\{\begin{aligned}
X_m^*&=\coprod_{\substack{i\in I_m\\\deleted{\unsure{d(i)}}}}X_i^*\\[1.5em]
V_{m,\deleted{\unsure{n}}\inserted{m+1}}^{\deleted{\unsure{*}}}&=\coprod_{\substack{(i,j)\in I\times I\\i\leq j,\ d(i)=m\\d(j)=m+1}}V_{i,j}\\[1.5em]
V_{m,m+1}^*&=\coprod_{\substack{i,j\text{ as}\\\text{above}}}V_{ij}^*
\end{aligned}\right.
\]
so that we have the diagrams
\diagram{\begin{tikzcd}[column sep=large,row sep=1em]
(26) & X_m^* \arrow[r] & V_{m,m+1}^{\deleted{\unsure{*}}} & V_{m,m+1}^* \arrow[l] \arrow[dd] \\
(\mathrm{Diag})_m & & & \\
& & & X_{m+1}^*
\end{tikzcd}}
fitting into a common diagram \unsure{\inserted{\unsure{\uline{which is none other than}|which \uline{is none other than}} (20)}}% doubt concerns placement after commun, not the inserted words.

infinite in principle, but which we
shall write in finite form by \deleted{\unsure{presuming|favouring}}
to simplify \unsure{\deleted{$I$ finite}|$I$ finite}, \inserted{$\unsure{d(I)|d(I)d}$ finite and $\subset\mathbb N$)} and by not writing the
\deleted{\unsure{empty} \unsure{components}} \inserted{\unsure{part|parts} of the diagram} \unsure{corresponding,|corresponding} to :
vertices occupied by empty spaces :
\margin{66}
